\section[Conclusão]{Conclusão}

Gostaria de iniciar com um prefácio: a AGES III era a etapa que eu mais aguardava cursar, tanto com ansiedade quanto com entusiasmo. No passado, eu já havia imaginado diversas vezes os princípios, padrões e até as tecnologias que gostaria de aplicar e discutir com o time. Além disso, eu finalmente teria a chance de aprofundar meus conhecimentos em infraestrutura na nuvem. Todavia, a experiência que tive acabou sendo o completo oposto de tudo o que imaginei.

A meu ver, nosso time não se entrosou e trabalhamos muito separados. Essa falta de entrosamento ficou evidente em praticamente qualquer discussão, seja sobre tecnologias, seja sobre a implementação das nossas tarefas. Já no início do projeto, inclusive na Sprint 0, deparei-me com colegas que me confrontaram agressivamente sobre decisões tanto de arquitetura quanto de projeto. Da minha parte, reconheço que me faltou uma “casca grossa” para lidar com essas situações; talvez, se eu tivesse agido de outra forma, o projeto pudesse ter alcançado um nível maior de qualidade e de entregas.

Também senti extrema falta de apoio por parte do time da AGES. Se não fosse pelo nosso trio da AGES III e por alguns colegas de outras disciplinas, não teríamos sequer descoberto que nós mesmos deveríamos criar o repositório da \textit{Wiki} — que, por sua vez, precisava conter os diagramas de \textit{deploy} necessários para a aprovação da infraestrutura —, algo que também não sabíamos onde encontrar dentro da área da AGES. Foi uma sequência de informações faltantes que prejudicou bastante nosso tempo de desenvolvimento e, consequentemente, nossa entrega aos clientes. Tenho a impressão de que, especificamente na nossa área de tecnologia, existe a expectativa de que já detenhamos, desde o início, todo o conhecimento necessário para cursar disciplinas mais específicas.

Refletindo sobre as lições que aprendi, vejo como ponto positivo o fato de que superar os empecilhos ao desenvolver a infraestrutura e ao desenhar a arquitetura de um projeto é, de fato, gratificante. Escolher tecnologias, guiar o time e auxiliar colegas a partir de conhecimentos prévios são experiências que levo para a minha carreira com bastante orgulho. Também percebi paralelos entre este projeto e minha vivência na AGES I, especialmente quanto ao desafio de trabalhar com um grupo não muito entrosado. De certa forma, nem todo mundo consegue trabalhar com todo mundo — e isso, por si só, é um desafio importante dentro da AGES.
